% GENERATED FILE. Edit canonical lessons, then run npm run generate:books.
% canonical-source-hash: fnv1a64:1c6a19bb2447834f
% canonical-lessons: ZH-C388-choukong, ZH-C388-yuehui, ZH-W388-yue-arrange, ZH-C388-shiyue, ZH-C388-tiqian, ZH-C388-shiqu

\chapter{Making a Date}
\label{ch:zh-making-a-date}
\hlchaptermodality{388}

\begin{chapteropening}
\textbf{By the end of this chapter:} \emph{I can talk about finding time, a date, missing an appointment, telling people early and time zones.}

Complete the last lesson of chapter 388: I can talk about finding time, a date, missing an appointment, telling people early and time zones.
\end{chapteropening}

\section[chōukòng]{chōukòng — to find time (for something)}

\label{lesson:ZH-C388-choukong}



\begin{quote}
Before the new one: say the Mandarin for to shop online, then the Mandarin for a budget.
\end{quote}

\subsection*{You'll want to know: chōukòng}
\emph{chōukòng} — \textquotedblleft{}to find time (for something)\textquotedblright{}.

\begin{quote}
\emph{yàoshi nǐ néng chōukòng lái kàn wǒ wǒ huì hěn gāoxìng} — If you can find time to come and see me, I'll be really pleased.
\end{quote}

Its characters come later; for now it is a word for the ear.

\begin{cousinweb}
\emph{chōukòng} is built from \emph{chōu} 'to pull out' and \emph{kòng} 'a gap', both inherited from Middle Chinese, \emph{trhjuw} and \emph{khuwngH} (Guangyun, in Baxter's transcription), which glosses the first 'to draw out' and the second 'a gap'.
\end{cousinweb}

\begin{grammarlens}[title={What you've built}]
One more word for saying what would happen.
\end{grammarlens}

\subsection*{Your turn}
\begin{itemize}
\raggedright
  \item \emph{Say it:} \emph{chōukòng}
  \item \emph{Say it:} \emph{chōukòng}, once more
  \item \emph{Say it:} \emph{yàoshi nǐ néng chōukòng lái kàn wǒ wǒ huì hěn gāoxìng} — If you can find time to come and see me, I'll be really pleased.
  \item \emph{Recall:} say the Mandarin for to shop online, then the Mandarin for a budget, then say \emph{chōukòng} again
\end{itemize}

\begin{tcolorbox}[breakable,colback=teal!4,colframe=teal!35!black,title={Before you move on}]
What does chōukòng mean? (\textquotedblleft{}To find time (for something)\textquotedblright{}.) Say it once more.
\end{tcolorbox}
\section[\texorpdfstring{\zh{约会} (yuēhuì)}{约会 (yuēhuì)}]{\zh{约会} (yuēhuì) — a romantic date; an appointment (to meet)}

\label{lesson:ZH-C388-yuehui}



\begin{quote}
Before the new one: say the Mandarin for a budget, then the Mandarin for to find time (for something).
\end{quote}

\subsection*{You'll want to know: \zh{约会}}
\textbf{\zh{约会}} — \emph{yuēhuì} — \textquotedblleft{}a romantic date; an appointment (to meet)\textquotedblright{}.

\begin{quote}
\textbf{\zh{今天晚上我有一个约会}} \emph{jīntiān wǎnshang wǒ yǒu yí ge yuēhuì} — I've got a date this evening.
\end{quote}

The new character, in pieces: \zh{约} — six strokes in two pieces: the silk piece on the left of \zh{给}, three strokes; on the right a short fall, a bar that turns down into a hook, and a dot inside. The next lesson puts it in your hand.

\begin{cousinweb}
\textbf{\zh{约会}} (\emph{yuēhuì}) is built from \zh{约} 'to bind' and \zh{会} 'to meet', both inherited from Middle Chinese, \emph{'jak} and \emph{hwajH} (Guangyun, in Baxter's transcription), which glosses the first 'to tie up' and the second 'to come together': a meeting agreed on beforehand.
\end{cousinweb}

\begin{grammarlens}[title={What you've built}]
One more word for saying what would happen.
\end{grammarlens}

\subsection*{Your turn}
\begin{itemize}
\raggedright
  \item \emph{Say it:} \emph{yuēhuì}
  \item \emph{Say it:} \emph{yuēhuì}, once more
  \item \emph{Say it:} \emph{jīntiān wǎnshang wǒ yǒu yí ge yuēhuì} — I've got a date this evening.
  \item \emph{Recall:} say the Mandarin for a budget, then the Mandarin for to find time (for something), then say \emph{yuēhuì} again
\end{itemize}

\begin{tcolorbox}[breakable,colback=teal!4,colframe=teal!35!black,title={Before you move on}]
What does \zh{约会} mean? (\textquotedblleft{}A romantic date; an appointment (to meet)\textquotedblright{}.) Say it once more.
\end{tcolorbox}
\section[\texorpdfstring{\zh{约} (yuē)}{约 (yuē)}]{\zh{约} — the silk piece of \zh{给}, then a hook around a dot}

\label{lesson:ZH-W388-yue-arrange}



\begin{quote}
The word first, then the character inside it.

\begin{itemize}
\raggedright
  \item \emph{Recall:} say \emph{yuēhuì} — \textbf{R1}, one lesson back, before it decays
\end{itemize}

\begin{quote}
\zh{给} \zh{会}
\end{quote}

\emph{Read it:} each one aloud, with nothing said first — \emph{gěi}, \emph{huì}
\end{quote}

\subsection*{Script — the silk piece of \zh{给}, then a hook around a dot}
\hlblockfigure{\detokenize{figures/ZH-W388-yue-arrange-filmstrip.pdf}}{How \zh{约} is written, stroke by stroke}

\begin{quote}
\zh{约}
\end{quote}

\textbf{Two pieces, left before right.}

\textbf{The left piece, the silk piece, three strokes,} as on the left of \zh{给}: two strokes that each fall to the left and turn to the right, one under the other, then a short stroke rising to the right.

\textbf{The right piece, three strokes.} A short stroke falling to the left. A bar that runs right and turns down into a long hook. Then a dot inside.

3 + 3 = six.

\subsection*{Your turn}
\begin{itemize}
\raggedright
  \item \emph{Write it:} \zh{约}, all six, with the model in front of you
  \item \emph{Cover:} write it again from memory, and count the pieces, not the strokes
\end{itemize}

\begin{tcolorbox}[breakable,colback=teal!4,colframe=teal!35!black,title={Before you move on}]
How many strokes in \zh{约}? (\textbf{Six}.) Where have you made the left piece before? (\textbf{On the left of \zh{给}}.) What comes last? (\textbf{The dot inside the hook}.)

Source: \href{https://raw.githubusercontent.com/chanind/hanzi-writer-data/68d10a4b21150cae5e1ebbd223eed289cf32d90c/data/\%E7\%BA\%A6.json}{Hanzi Writer Data \zh{约}, pinned stroke paths and medians}.
\end{tcolorbox}
\section[shīyuē]{shīyuē — to miss an appointment}

\label{lesson:ZH-C388-shiyue}



\begin{quote}
Before the new one: say the Mandarin for to find time (for something), then the Mandarin for a romantic date.
\end{quote}

\subsection*{You'll want to know: shīyuē}
\emph{shīyuē} — \textquotedblleft{}to miss an appointment\textquotedblright{}.

\begin{quote}
\emph{rúguǒ nǐ shīyuē tā huì hěn shēngqì} — If you don't turn up, he'll be angry.
\end{quote}

Its characters come later; for now it is a word for the ear.

\begin{cousinweb}
\emph{shīyuē} is built from \emph{shī} 'to lose' and \zh{约} 'an agreement', both inherited from Middle Chinese, \emph{syit} and \emph{'jak} (Guangyun, in Baxter's transcription), which glosses the first 'to let slip'.
\end{cousinweb}

\begin{grammarlens}[title={What you've built}]
One more word for saying what would happen.
\end{grammarlens}

\subsection*{Your turn}
\begin{itemize}
\raggedright
  \item \emph{Say it:} \emph{shīyuē}
  \item \emph{Say it:} \emph{shīyuē}, once more
  \item \emph{Say it:} \emph{rúguǒ nǐ shīyuē tā huì hěn shēngqì} — If you don't turn up, he'll be angry.
  \item \emph{Recall:} say the Mandarin for to find time (for something), then the Mandarin for a romantic date, then say \emph{shīyuē} again
\end{itemize}

\begin{tcolorbox}[breakable,colback=teal!4,colframe=teal!35!black,title={Before you move on}]
What does shīyuē mean? (\textquotedblleft{}To miss an appointment\textquotedblright{}.) Say it once more.
\end{tcolorbox}
\section[tíqián]{tíqián — to do ahead of time; to bring forward (a date)}

\label{lesson:ZH-C388-tiqian}



\begin{quote}
Before the new one: say the Mandarin for a romantic date, then the Mandarin for to miss an appointment.
\end{quote}

\subsection*{You'll want to know: tíqián}
\emph{tíqián} — \textquotedblleft{}to do ahead of time; to bring forward (a date)\textquotedblright{}.

\begin{quote}
\emph{yàoshi nǐ bù néng lái zuìhǎo tíqián gàosu wǒ} — If you can't come, you'd better tell me in good time.
\end{quote}

Its characters come later; for now it is a word for the ear.

\begin{cousinweb}
\emph{tíqián} is built from \emph{tí} 'to carry by the handle' and \zh{前} 'before', both inherited from Middle Chinese, \emph{dej} and \emph{dzen} (Guangyun, in Baxter's transcription), which glosses the second 'first': lifting a time forward.
\end{cousinweb}

\begin{grammarlens}[title={What you've built}]
One more word for saying what would happen.
\end{grammarlens}

\subsection*{Your turn}
\begin{itemize}
\raggedright
  \item \emph{Say it:} \emph{tíqián}
  \item \emph{Say it:} \emph{tíqián}, once more
  \item \emph{Say it:} \emph{yàoshi nǐ bù néng lái zuìhǎo tíqián gàosu wǒ} — If you can't come, you'd better tell me in good time.
  \item \emph{Recall:} say the Mandarin for a romantic date, then the Mandarin for to miss an appointment, then say \emph{tíqián} again
\end{itemize}

\begin{tcolorbox}[breakable,colback=teal!4,colframe=teal!35!black,title={Before you move on}]
What does tíqián mean? (\textquotedblleft{}To do ahead of time; to bring forward (a date)\textquotedblright{}.) Say it once more.
\end{tcolorbox}
\section[\texorpdfstring{\zh{时区} (shíqū)}{时区 (shíqū)}]{\zh{时区} (shíqū) — time zone}

\label{lesson:ZH-C388-shiqu}



\begin{quote}
Before the new one: say the Mandarin for to miss an appointment, then the Mandarin for to do ahead of time.
\end{quote}

\subsection*{You'll want to know: \zh{时区}}
\textbf{\zh{时区}} — \emph{shíqū} — \textquotedblleft{}time zone\textquotedblright{}.

\begin{quote}
\textbf{\zh{我们不在一个时区所以要先约好时间}} \emph{wǒmen bú zài yí ge shíqū suǒyǐ yào xiān yuē hǎo shíjiān} — We're not in the same time zone, so we need to fix a time first.
\end{quote}

\begin{cousinweb}
\textbf{\zh{时区}} (\emph{shíqū}) is built from \zh{时} 'time' and \zh{区} 'an area', both inherited from Middle Chinese, \emph{dzyi} and \emph{khju} (Guangyun, in Baxter's transcription), which glosses the first 'an hour of the day'. The word is modern.
\end{cousinweb}

\begin{grammarlens}[title={What you've built}]
One more word for saying what would happen.
\end{grammarlens}

\subsection*{Your turn}
\begin{itemize}
\raggedright
  \item \emph{Say it:} \emph{shíqū}
  \item \emph{Say it:} \emph{shíqū}, once more
  \item \emph{Say it:} \emph{wǒmen bú zài yí ge shíqū suǒyǐ yào xiān yuē hǎo shíjiān} — We're not in the same time zone, so we need to fix a time first.
  \item \emph{Recall:} say the Mandarin for to miss an appointment, then the Mandarin for to do ahead of time, then say \emph{shíqū} again
\end{itemize}

\begin{tcolorbox}[breakable,colback=teal!4,colframe=teal!35!black,title={Before you move on}]
What does \zh{时区} mean? (\textquotedblleft{}Time zone\textquotedblright{}.) Say it once more.
\end{tcolorbox}
